% =====================================================
% CHAPITRE 4 : PRÉPARATION DES DONNÉES (CRISP-DM Phase 3)
% =====================================================
\chapter{Préparation des données}

% ---------------------------------------------------
\section*{Introduction}
\addcontentsline{toc}{section}{Introduction}
% ---------------------------------------------------

Préparer les données, c'est poser les fondations sur lesquelles reposent tous les modèles
à venir. J'ai rapidement réalisé que les décisions prises ici auraient des répercussions
directes sur la validité de l'évaluation finale. L'enjeu central n'était pas seulement
technique, mais méthodologique : comment extraire des signaux pertinents des variables
brutes, tout en évitant d'introduire sans le vouloir de l'information du futur dans
l'apprentissage --- ce que l'on appelle le \textit{data leakage} ? Ce chapitre détaille
chaque choix effectué et les justifications empiriques qui les sous-tendent, telles
qu'elles ressortent des notebooks \texttt{01\_data\_understanding.ipynb} et
\texttt{02\_data\_preparation.ipynb}.

% ---------------------------------------------------
\section{Pipeline de préparation}
% ---------------------------------------------------

\subsection{Vue d'ensemble}

Le pipeline de préparation suit les étapes définies dans le module
\texttt{src/preprocessing/preprocessing.py}. Ces étapes sont appliquées dans un ordre
strict pour garantir l'étanchéité des partitions et l'absence de fuite d'information :

\begin{enumerate}
  \item Suppression des variables à risque de leakage
  \item Ingénierie de variables (\textit{feature engineering})
  \item Division stratifiée du corpus (train / validation / test)
  \item Normalisation (StandardScaler, \textit{fit} sur train uniquement)
  \item Gestion du déséquilibre (pondération et SMOTE)
\end{enumerate}

\subsection{Suppression des variables à risque de leakage}

La première décision consiste à identifier et éliminer les variables qui introduiraient de
l'information future ou qui révèleraient directement la classe cible. Ces suppressions sont
motivées empiriquement par l'analyse exploratoire :

\begin{table}[H]
\centering
\caption{Variables supprimées et justifications}
\label{tab:cols_dropped}
\begin{tabular}{p{4cm}p{8cm}}
\toprule
\textbf{Variable} & \textbf{Justification} \\
\midrule
\texttt{nameOrig}, \texttt{nameDest}
  & Identifiants nominaux non encodables sans risque de sur-apprentissage \\[0.2cm]
\texttt{oldbalanceOrg}, \texttt{newbalanceOrig}
  & Corrélation directe avec la fraude ; fuite d'information avérée \\[0.2cm]
\texttt{oldbalanceDest}, \texttt{newbalanceDest}
  & Idem ; révèlent les patterns frauduleux de manière trop évidente \\[0.2cm]
\texttt{isFlaggedFraud}
  & Label proxy du système existant ; ne doit pas alimenter l'apprentissage \\
\bottomrule
\end{tabular}
\end{table}

% ---------------------------------------------------
\section{Ingénierie de variables}
% ---------------------------------------------------

\subsection{Variables temporelles}

La variable \texttt{step} (pas de temps en heures) est décomposée en trois variables
cycliques qui capturent des patterns à différentes granularités temporelles :

\begin{itemize}
  \item \texttt{hour} : heure dans la journée (0--23), calculée par
    $\texttt{step} \bmod 24$ ;
  \item \texttt{day} : jour de simulation (0--30), calculé par
    $\lfloor \texttt{step} / 24 \rfloor \bmod 31$ ;
  \item \texttt{week} : semaine de simulation (0--4), calculée par
    $\lfloor \texttt{step} / 168 \rfloor$.
\end{itemize}

\subsection{Variables comportementales et de risque}

Cinq variables dérivées sont construites à partir des observations de l'EDA :

\begin{table}[H]
\centering
\caption{Variables dérivées construites lors du feature engineering}
\label{tab:features_derivees}
\begin{tabular}{p{3.5cm}p{4.5cm}p{4cm}}
\toprule
\textbf{Variable} & \textbf{Formule / Logique} & \textbf{Signal mesuré} \\
\midrule
\texttt{high\_risk\_hour} &
  $1$ si \texttt{hour} $\in \{0,...,9, 23\}$ &
  Ratio fraude $10{,}63\times$ supérieur \\[0.2cm]
\texttt{is\_transfer\_or\_cashout} &
  $1$ si \texttt{type} $\in$ \{TRANSFER, CASH\_OUT\} &
  Seuls types avec fraudes réelles \\[0.2cm]
\texttt{balance\_diff\_orig} &
  \texttt{oldbalanceOrg} $-$ \texttt{newbalanceOrig} &
  Corrélation 0{,}3662 avec \texttt{isFraud} \\[0.2cm]
\texttt{dest\_zero\_balance} &
  $1$ si \texttt{oldbalanceDest} $= 0$ &
  Taux fraude 2{,}52\,\% vs 0{,}13\,\% global \\[0.2cm]
\texttt{log\_amount} &
  $\log(1 + \texttt{amount})$ &
  Corrige skewness $= 30{,}80$ \\
\bottomrule
\end{tabular}
\end{table}

Chacune de ces variables est justifiée empiriquement par l'EDA --- aucun choix n'est
arbitraire.

\subsection{Encodage one-hot}

La variable catégorielle \texttt{type} (5 modalités) est encodée par
\textbf{one-hot encoding} (\textit{drop\_first=False}), produisant les colonnes
\texttt{type\_CASH\_IN}, \texttt{type\_CASH\_OUT}, \texttt{type\_DEBIT},
\texttt{type\_PAYMENT} et \texttt{type\_TRANSFER}. L'absence de suppression d'une
colonne (\textit{drop\_first=False}) est volontaire : avec un arbre de décision ou un
autoencoder, toutes les modalités sont nécessaires.

\subsection{Ensemble final des variables}

Après feature engineering, le jeu de données comprend \textbf{14 variables} :

\begin{table}[H]
\centering
\caption{Ensemble final des 14 variables du modèle}
\label{tab:final_features}
\begin{tabular}{p{2.5cm}p{6cm}p{3cm}}
\toprule
\textbf{Catégorie} & \textbf{Variables} & \textbf{Traitement} \\
\midrule
Temporelles  & \texttt{step}, \texttt{hour}, \texttt{day}, \texttt{week} & StandardScaler \\
Montant      & \texttt{log\_amount} & StandardScaler \\
Comportement & \texttt{balance\_diff\_orig} & StandardScaler \\
Binaires     & \texttt{high\_risk\_hour}, \texttt{is\_transfer\_or\_cashout},
               \texttt{dest\_zero\_balance} & --- \\
Type (OHE)   & \texttt{type\_CASH\_IN}, \texttt{type\_CASH\_OUT}, \texttt{type\_DEBIT},
               \texttt{type\_PAYMENT}, \texttt{type\_TRANSFER} & --- \\
\bottomrule
\end{tabular}
\end{table}

% ---------------------------------------------------
\section{Division et normalisation}
% ---------------------------------------------------

\subsection{Division stratifiée}

Le corpus est divisé en trois sous-ensembles de manière \textbf{stratifiée} sur la variable
cible, afin de conserver la proportion de fraudes dans chaque partition --- essentiel
lorsque les classes sont aussi déséquilibrées.

\begin{table}[H]
\centering
\caption{Division stratifiée du corpus}
\label{tab:splits}
\begin{tabular}{lrrl}
\toprule
\textbf{Ensemble} & \textbf{Taille} & \textbf{Fraudes} & \textbf{Rôle} \\
\midrule
Entraînement & 139~999 & 181 & Apprentissage des modèles \\
Validation   & 30~000  & 38  & Réglage des hyperparamètres et seuils \\
Test         & 30~001  & 39  & Évaluation finale (jamais consulté avant) \\
\bottomrule
\end{tabular}
\end{table}

\textbf{Principe d'étanchéité :} L'ensemble de test n'est jamais consulté avant
l'évaluation finale, conformément aux bonnes pratiques anti-\textit{data snooping}.
Ce principe a été respecté tout au long du projet.

\subsection{Normalisation}

Le \textbf{StandardScaler} est ajusté (\textit{fit}) exclusivement sur l'ensemble
d'entraînement, puis appliqué (\textit{transform}) sur les trois partitions. Cette
discipline stricte garantit l'absence de leakage statistique : aucune information sur la
distribution du test ne s'infiltre dans les transformations appliquées à l'entraînement.

% ---------------------------------------------------
\section{Gestion du déséquilibre}
% ---------------------------------------------------

\subsection{Pondération des classes}

Pour les modèles supervisés, des poids de classes sont calculés automatiquement
(formule \textit{sklearn balanced}) :

\[
w_c = \frac{n_{\text{total}}}{n_{\text{classes}} \times n_c}
\]

Avec 139~999 transactions d'entraînement dont 181 fraudes, les poids obtenus sont
approximativement $w_0 \approx 0{,}50$ et $w_1 \approx 386{,}7$. Cela revient à donner
à chaque transaction frauduleuse un poids environ 386 fois supérieur à une transaction
légitime lors du calcul de la fonction de perte --- compensant ainsi le déséquilibre du
corpus.

\subsection{SMOTE}

Pour les variantes avec rééchantillonnage, \textbf{SMOTE} (\textit{Synthetic Minority
Over-sampling Technique}) génère des exemples synthétiques de la classe minoritaire par
interpolation dans l'espace des features. Le principe : pour chaque transaction frauduleuse,
SMOTE choisit aléatoirement $k$ voisins les plus proches dans l'espace des features, puis
crée de nouveaux exemples sur les segments reliant la transaction à ses voisins.

Il est appliqué \textbf{uniquement sur l'ensemble d'entraînement} (jamais sur validation
ou test), afin de ne pas contaminer l'évaluation avec des exemples synthétiques.

\subsection{Ensemble pour l'autoencoder}

Pour l'autoencoder non supervisé, un sous-ensemble \texttt{X\_train\_normal} est extrait :
il contient \textbf{uniquement les 139~818 transactions légitimes} de l'entraînement
(0 fraude). Cela permet au modèle d'apprendre exclusivement la distribution normale ---
principe fondateur de la détection d'anomalies par reconstruction.

% ---------------------------------------------------
\section*{Conclusion}
\addcontentsline{toc}{section}{Conclusion}
% ---------------------------------------------------

La préparation des données produit un jeu de \textbf{14 variables représentatives} des
patterns de fraude, en respectant strictement l'étanchéité entre les partitions. Les
7 variables dérivées ont chacune été validées par une mesure empirique issue de l'EDA ---
aucune n'est arbitraire. La discipline anti-leakage (StandardScaler fit sur train seulement,
SMOTE appliqué uniquement sur le train, test jamais consulté) garantit que l'évaluation
finale reflète les véritables performances sur des données inédites. Ce socle solide permet
d'entraîner des modèles fiables et comparables, décrits dans le chapitre suivant.
